\chapter{Two-Timescale Dynamics}
\label{ch:two-timescales}

This chapter asks what the collective free energy of \cref{ch:elbo} supplies for
separating fast inference of beliefs from slower evolution of models, priors, and
parents: descent for every positive-definite block mobility, local in a product family
(\cref{prop:obs-local-block-descent}); a reduced slow flow conjugate to its
tangent lift on the graph of the conditional optimum
(\cref{lem:obs-envelope-slow-flow,prop:obs-rate-separated-flow}), the exact route of
\cref{ch:relational-inference} for information-discarding maps; and a Markov parent
history under instantaneous within-parent refresh
(\cref{thm:rg-memory-prepared-parent}), whose failure under finite refresh
(\cref{rem:rg-memory-partial-refresh}) is measured in \cref{ch:agent-network-rg}. The
slow block (\cref{def:obs-fast-slow-blocks}) moves by accepted M-coordinates
(\cref{def:elbo-generalized-m-acceptance}) or by the rate-separated flow; the formation
family of \cref{ch:finite-depth-formation} realizes one scalar slow parameter
\eqref{eq:rg-formation-own-law-dynamics}, not slow per-agent model laws, while on
the four-cycle of that chapter, whose link translations this chapter writes
$\check\omega$, slow model means are a supplied rate, and the model
channel lags the belief channel only when the links are slow as well, at a rate of
order at most the model rate
(\cref{prop:obs-slow-model-link-dichotomy}, the held-sheet instance of
\cref{open:obs-slow-manifold}). The finite
model of the last section uses fine states $\xi\in\Xi$, parent values $z\in\mathsf Z$,
a parent map $\operatorname{par}$, a prior $p^0$, record laws $\psi_\xi$, a refresh
fraction $\omega$, and sans-serif kernels acting on row vectors from the right.

\paragraph{Scope.}
The collective free energy does not choose the fast--slow split, the rate ratio, or
the block mobilities, and a rate ratio is an anisotropic mobility that can change the
selected orbit (\cref{prop:hist-scalar-mobility-orbit}). The chapter does not claim
that the rate-separated flow tracks the reduced flow for a recognition family of the
program beyond the held matching sheets of
\cref{prop:obs-slow-model-link-dichotomy}, that the sampled dynamics with source
jumps track it, that a sampled episode with source jumps at the temperature of
\cref{def:rg-formation-cycle-experiment} keeps a matched pair belief-coherent and
model-incoherent on an interval of order $1/\varepsilon_{\mathrm m}$ with probability
bounded below, that the discrete alternating scheme converges, that the parent sweep
converges at a rate uniform in the number of fine states or in the smallest posterior
weight, or that a static parent datum has closed dynamics
(\cref{open:obs-slow-manifold,open:elbo-alternating-convergence,open:rg-pointwise-parent-dynamics}).
\status{NOT-CLAIMED}

\section{Block descent at fixed parameter}
\label{sec:local-natural-gradient}

\propositionheading{Block descent of the collective free energy}{prop:obs-local-block-descent}
Let $\lambda\mapsto Q_\lambda$ be a family of recognition laws on an open domain
$\Lambda=\prod_{i\in V}\Lambda_i$ with block coordinates $\lambda=(\lambda_i)_i$,
such that $F(\lambda):=\Fenergy^{\mathrm{ext}}_o(Q_\lambda)$ is real and $C^1$ on
$\Lambda$, and let $\mathsf G_i(\lambda)$ be continuous fields of symmetric
positive-definite block metrics, with positive block rates absorbed into the
metrics and blocks of zero rate held fixed. Let $\lambda(t)$ solve
\begin{equation}
\dot\lambda_i=-\mathsf G_i^{-1}\nabla_{\lambda_i}F
\label{eq:obs-local-natural-gradient}
\end{equation}
for the moving blocks $i\in V'\subseteq V$ and remain in $\Lambda$.
\begin{enumerate}
\item[(i)] Along the solution
\begin{equation}
\frac{d}{dt}F=-\sum_{i\in V'}(\nabla_{\lambda_i}F)^\top\mathsf G_i^{-1}(\nabla_{\lambda_i}F)\leq0.
\label{eq:obs-global-dissipation}
\end{equation}
\item[(ii)] If the family is a regular product family
$Q_\lambda=\bigotimes_iQ_{i,\lambda_i}$ with
$\E_{Q_{-i}}\lvert\log Z_{\{i\}}(Y_{-i})\rvert<\infty$ for every $i$ and $\lambda$,
then $\nabla_{\lambda_i}F=\nabla_{\lambda_i}\overline\Fenergy_i$ for the
outside-averaged local free energy \eqref{eq:obs-outside-averaged-local-vfe}, so
an agent that descends its own local free energy moves along a block of a
descent flow of the collective one. The Fisher metric of such a family is block
diagonal, so with $\mathsf G_i$ equal to the block Fisher metrics the flow is the
Fisher natural gradient of $F$.
\item[(iii)] Outside product families moving $\lambda_i$ can change $Q_{-i}$,
and the identification in (ii) needs the conditional lift of
\cref{hyp:hist-exact-vfe-lift}, whose outside-averaged objective
\eqref{eq:hist-outside-averaged-local-vfe} has the differential of the collective
free energy restricted to fixed outside marginal. The Fisher metric then need not
be block diagonal: for the bivariate normal location family with unit variances
and correlation $\mathrm c_2$, $0<\lvert\mathrm c_2\rvert<1$, its off-diagonal entry
is $-\mathrm c_2/(1-\mathrm c_2^2)$.
\end{enumerate}
\status{ESTABLISHED}

\noindent\emph{Proof.}
Part (i) is the chain rule for $F$ along the solution. For (ii), a change of
$\lambda_i$ leaves $Q_{-i}$ fixed, and \cref{thm:obs-local-global-potential} with
$B=\{i\}$, whose finiteness hypotheses hold because $F$ is real, shows that $F$
and $\overline\Fenergy_i$ differ by a quantity independent of $\lambda_i$. Scores
of distinct factors of a product law are independent with mean zero, so their
cross covariances vanish. For (iii), the Fisher metric of the location family is
the inverse covariance matrix. $\square$

\remarkheading{Simultaneous block updates}{rem:obs-simultaneous-block-updates}
The flow above, \eqref{eq:obs-damped-block-update}, and
\cref{cor:obs-potential-game} move one block at a time or move continuously; a
discrete schedule that replaces several blocks at once is covered only in the
following cases. (a) Let $B_1,\dots,B_k$ be disjoint blocks with union $U$ such
that no record scope meets two of them and the baseline makes
$Y_{B_1},\dots,Y_{B_k}$ conditionally independent given $Y_{-U}$, as a product
baseline does; for singleton blocks under a product baseline this is an
independent set of the record hypergraph. For a law
$Q=Q_{-U}\bigotimes_lr_{B_l}(\cdot\given y_{-U})$ with
$\Fenergy^{\mathrm{ext}}_o(Q)<\infty$, replacing every $r_{B_l}$ at once by a damped
kernel of \cref{prop:elbo-exact-e-coordinate}(iv) gives the law reached by the
sequential sweep in any order, so it does not increase $\Fenergy^{\mathrm{ext}}_o$;
and under the hypotheses of \cref{cor:obs-potential-game}, changing the strategies
of the agents of such a set at once changes $\Fenergy^{\mathrm{ext}}_o$ by the sum
of their unilateral changes. A chromatic schedule that updates the classes of a
partition of the agents into such sets in turn therefore descends whenever each
unilateral change does. (b) In \cref{cor:obs-potential-game} let the strategy set of
agent $i$ be a regular minimal exponential family
$q_{i,\eta}(dy_i)=\exp\{\langle\eta,\mathsf T_i(y_i)\rangle-\mathsf A_i(\eta)\}\nu_i(dy_i)$,
$\eta\in\mathsf N_i$, with mean parameter $\mu=\nabla\mathsf A_i(\eta)$, and let it
be conditionally conjugate: with the other strategies held,
$\overline\Fenergy_i(q_{i,\eta};r_{-i})=\KL(q_{i,\eta}\Vert q_{i,\eta^*})+\mathrm c_i$
for some $\eta^*\in\mathsf N_i$ and $\mathrm c_i$ independent of $\eta$, so that
$q_{i,\eta^*}$ is the best response. Damping in mean coordinates, to the member with
mean $(1-\mathrm w)\mu+\mathrm w\mu^*$ when that mean belongs to the family, which it
always does for Gaussian families, satisfies the two inequalities of
\eqref{eq:obs-damped-block-update} with the best response in place of the local
posterior. Damping in natural coordinates, to
$\eta_{\mathrm w}=(1-\mathrm w)\eta+\mathrm w\eta^*$, does not increase
$\Fenergy^{\mathrm{ext}}_o$ either, but it need not satisfy the first inequality:
for one agent with a Poisson latent, posterior $\operatorname{Poi}(1)$, natural
parameter the logarithm of the rate, $\eta=-4$, $\eta^*=0$ and $\mathrm w=1/2$, the
relative entropy to the posterior is $1-3e^{-2}=0.594\ldots$, above
$(1-5e^{-4})/2=0.454\ldots$. (c) In the conjugate case the natural gradient of
$\Fenergy^{\mathrm{ext}}_o$ in $\eta_i$ for the Fisher block metric is
$\eta_i-\eta_i^*$, so simultaneous natural-coordinate damping of all agents with
weight $\mathrm w$ is the explicit Euler step of size $\mathrm w$ of
\eqref{eq:obs-local-natural-gradient}, and $\mathrm w=1$ is the simultaneous best
response. For blocks that share records it can increase the free energy. Take three
agents with scalar latents, baseline $\mathcal N(0,5/2)$ for each, and one record
$o_a\sim\mathcal N(y_1+y_2+y_3,5/3)$, so the posterior is Gaussian with precision
$\Lambda=\tfrac25I+\tfrac35\mathbf 1\mathbf 1^\top$ and some mean $\bar y^\star$, and
take product Gaussian strategies. At unit variances and means $\bar y^\star+\mathbf 1$,
\begin{equation}
\KL(Q\Vert\boldsymbol\Pi_{\theta,V,o,X})=\frac{33}{10}+\frac12\log\frac{125}{44},
\label{eq:obs-simultaneous-witness}
\end{equation}
and the simultaneous best responses keep unit variances and move the means to
$\bar y^\star-\tfrac65\mathbf 1$, where the relative entropy is
$\tfrac{594}{125}+\tfrac12\log\tfrac{125}{44}$; each repetition multiplies the
first term by $36/25$, whereas the sequential sweep never increases the free energy.
The update of weight $\mathrm w$ increases it along $\mathbf 1$ exactly when
$\mathrm w>10/11=2/\mathrm c_{\mathsf G}$, where $\mathrm c_{\mathsf G}=11/5$ is the
largest eigenvalue of the Hessian $\Lambda$ relative to the Fisher metric $I$ of the
means at unit variance, which is the threshold of
\eqref{eq:hist-discrete-descent-bound}. A simultaneous step of blocks that share
records is therefore covered only by a step bound of that kind or by a step
accepted on the same functional, as after
\cref{prop:elbo-finite-step-nonmonotonicity}. \status{ESTABLISHED}

\noindent\emph{Proof.}
For (a), the incident likelihood $L_{B_l,o}$ of \eqref{eq:obs-block-incident-likelihood}
reads no coordinate of another $B_m$, and conditional independence gives the baseline
conditional of $Y_{B_l}$ given $Y_{-B_l}$ a version that reads only $y_{-U}$; by
\eqref{eq:obs-local-posterior} and \eqref{eq:obs-local-vfe} the local evidence,
posterior, and free energy of $B_l$ then read only $y_{-U}$. Under $Q$ the conditional
of $Y_{B_1}$ given $Y_{-B_1}$ is $r_{B_1}(\cdot\given y_{-U})$, so
\cref{prop:elbo-exact-e-coordinate}(iv) applies with $B=B_1$, and the damped kernel
again reads only $y_{-U}$, so the new law has the same form and finite free energy;
induction over the blocks, in any order, ends at
$Q_{-U}\bigotimes_lr^{\mathrm w_l}_{B_l}$. For product laws, the integrand of
\eqref{eq:obs-outside-averaged-local-vfe} for agent $i$ does not read $y_j$ for
another agent $j$ of the set, so $\overline\Fenergy_i$ does not depend on $r_j$;
changing the strategies one agent at a time, \cref{cor:obs-potential-game} gives
each change of $\Fenergy^{\mathrm{ext}}_o$ as the change of $\overline\Fenergy_i$ at
the original strategies of the set, and the changes add. For (b), $\mathsf A_i$ is
finite near $\eta$ and $\eta^*$, so $\mathsf T_i$ is integrable under both members,
with means $\mu=\nabla\mathsf A_i(\eta)$ and $\mu^*$, and the Fisher metric of the
natural chart is $\nabla^2\mathsf A_i$
\citep[Theorems~7.1, 7.2 and~8.1]{BarndorffNielsen1978}. By
\cref{cor:obs-potential-game}, $\Fenergy^{\mathrm{ext}}_o$ changes by the change of
$\KL(\cdot\Vert q_{i,\eta^*})$. The mixture
$\mathrm m=(1-\mathrm w)q_{i,\eta}+\mathrm wq_{i,\eta^*}$ has mean
$\mu_{\mathrm w}=(1-\mathrm w)\mu+\mathrm w\mu^*$; if $q_{i,\eta'_{\mathrm w}}$ has
that mean, $\log(q_{i,\eta'_{\mathrm w}}/q_{i,\eta^*})$ is affine in $\mathsf T_i$, so
its expectations under $\mathrm m$ and under $q_{i,\eta'_{\mathrm w}}$ agree, and
\[
(1-\mathrm w)\KL(q_{i,\eta}\Vert q_{i,\eta^*})\geq\KL(\mathrm m\Vert q_{i,\eta^*})
=\KL(\mathrm m\Vert q_{i,\eta'_{\mathrm w}})+\KL(q_{i,\eta'_{\mathrm w}}\Vert q_{i,\eta^*})
\geq\KL(q_{i,\eta'_{\mathrm w}}\Vert q_{i,\eta^*}),
\]
the first inequality by convexity of relative entropy in its first argument. This is
the first inequality of \eqref{eq:obs-damped-block-update} after adding the
constant, and the second follows because the best response minimizes. For a
Gaussian family the means $(\mu,\mathrm M_2)$ of $(Y,YY^\top)$ belong to the family exactly when
$\mathrm M_2-\mu\mu^\top\succ0$, that is, when the matrix with blocks $1$, $\mu^\top$, $\mu$, $\mathrm M_2$ is
positive definite, a convex condition. With $\Delta\eta=\eta-\eta^*$ and
$\mathrm f(t)=\KL(q_{i,\eta^*+t\Delta\eta}\Vert q_{i,\eta^*})
=\mathsf A_i(\eta^*)-\mathsf A_i(\eta^*+t\Delta\eta)+t\langle\nabla\mathsf A_i(\eta^*+t\Delta\eta),\Delta\eta\rangle$,
one has $\mathrm f'(t)=t\Delta\eta^\top\nabla^2\mathsf A_i(\eta^*+t\Delta\eta)\Delta\eta\geq0$ on the segment,
which lies in the convex domain $\mathsf N_i$, and $\eta_{\mathrm w}$ is the point
$t=1-\mathrm w$. In the Poisson family $\mathsf A(\eta)=e^\eta$ and
$\mathrm f$ at $\eta^*=0$ is $1-(1-\eta)e^{\eta}$ evaluated at $\eta=-2$ and
$\eta=-4$. For (c), differentiating the Bregman form gives
$\nabla_{\eta}\KL(q_{i,\eta}\Vert q_{i,\eta^*})=\nabla^2\mathsf A_i(\eta)(\eta-\eta^*)$,
which by \cref{prop:obs-local-block-descent}(ii) is $\nabla_{\eta_i}\Fenergy^{\mathrm{ext}}_o$;
multiplying by the inverse Fisher metric gives $\eta_i-\eta_i^*$, and the Euler step
of size $\mathrm w$ is $\eta_{\mathrm w}$. In the witness the record adds precision
$\tfrac35\mathbf 1\mathbf 1^\top$ to the baseline precision $\tfrac25I$, so
$\Lambda$ has unit diagonal and eigenvalues $\tfrac25,\tfrac25,\tfrac{11}5$ with
$\det\Lambda=44/125$. For a product law with variances $\mathrm v_i$ and mean error $\boldsymbol\delta$,
$\KL(Q\Vert\boldsymbol\Pi_{\theta,V,o,X})
=\frac12[\sum_i\mathrm v_i+\boldsymbol\delta^\top\Lambda\boldsymbol\delta-3-\log\det\Lambda-\sum_i\log\mathrm v_i]$, whose
best response for agent $i$ is $\mathrm v_i=1$ and $\delta_i=-\sum_{j\ne i}\Lambda_{ij}\delta_j$; at
$\boldsymbol\delta=\mathbf 1$ this is $\delta_i=-\tfrac65$, and
$\mathbf 1^\top\Lambda\mathbf 1=\tfrac{33}5$ gives both values. The weighted update
is $\boldsymbol\delta\mapsto\boldsymbol\delta-\mathrm w\Lambda\boldsymbol\delta$, which multiplies the component along
$\mathbf 1$ by $1-\tfrac{11}5\mathrm w$. Each sequential best response is a
unilateral improvement of \cref{cor:obs-potential-game}. $\square$

\section{Fast recognition and slow model coordinates}
\label{sec:obs-fast-slow}

\definitionheading{Fast and slow coordinate blocks}{def:obs-fast-slow-blocks}
Let $\mathcal U_{\rm f}\subseteq\R^{n_{\rm f}}$ and
$\mathcal U_{\rm s}\subseteq\R^{n_{\rm s}}$ be open, and write
$s=(s_{\rm m},\theta)\in\mathcal U_{\rm s}$. A fast--slow coordinatization is a
map $(f,s)\mapsto(Q_{f,s_{\rm m}},\theta)$ into population recognition laws and
generative parameters, together with the real-valued objective
\begin{equation}
F(f,s):=\Fenergy^{\mathrm{ext}}_o(Q_{f,s_{\rm m}};\theta)
=-\log p_{\theta,V}(o\given X)
+\KL\bigl(Q_{f,s_{\rm m}}\Vert\boldsymbol\Pi_{\theta,V,o,X}\bigr),
\label{eq:obs-fast-slow-objective}
\end{equation}
in which the generative argument of \eqref{eq:obs-collective-vfe} is explicit. The
fast coordinate $f$ carries the belief-channel recognition coordinates and, when
present, the source-label rows of \cref{sec:local-attention-elbo}; the slow
coordinate carries the model-channel recognition coordinates $s_{\rm m}$, which move
a recognition law, and the generative parameter $\theta$, which moves the fixed joint
between evaluations and never inside one, so a composite $\theta(Q)$ inside one
evaluation of \eqref{eq:obs-fast-slow-objective} stays excluded by
\cref{req:gen-typing-prohibition}. Because \eqref{eq:obs-fast-slow-objective}
evaluates one observation $o$ at every $\theta$ in a continuum, the coordinatization
also declares, as model data in the sense of
\cref{hyp:prob-regular-conditional-convention}, a version
$(\theta,o,y)\mapsto p_{\theta,V}(o,y\given X)$ jointly measurable in its arguments,
with evidence $0<p_{\theta,V}(o\given X)<\infty$ at the observed $o$ for every
$\theta$ occurring in $\mathcal U_{\rm s}$ and posterior slices computed from that
version by \eqref{eq:prob-rcp-density}; real values and differentiability of $F$
in $\theta$ are statements about this version. The assignment of coordinates to
the two blocks is a declaration.

\lemmaheading{Envelope identity and reduced slow descent}{lem:obs-envelope-slow-flow}
Let $F\in C^1(\mathcal U_{\rm f}\times\mathcal U_{\rm s})$, let
$\mathcal U_0\subseteq\mathcal U_{\rm s}$ be open, and let
$f^\star:\mathcal U_0\to\mathcal U_{\rm f}$ be a $C^1$ map with
$F(f^\star(s),s)=\min_{f\in\mathcal U_{\rm f}}F(f,s)$ for every $s\in\mathcal U_0$.
The reduced functional
\begin{equation}
F_{\rm red}(s):=\min_{f\in\mathcal U_{\rm f}}F(f,s)=F(f^\star(s),s)
\label{eq:obs-reduced-functional}
\end{equation}
satisfies the envelope identity
\begin{equation}
dF_{\rm red}(s)=\partial_sF(f^\star(s),s),\qquad s\in\mathcal U_0.
\label{eq:obs-envelope-identity}
\end{equation}
For every continuous field $\mathsf G_{\rm s}$ of symmetric positive-definite
metrics on $\mathcal U_{\rm s}$, the reduced slow flow
\begin{equation}
\dot s=-\mathsf G_{\rm s}(s)^{-1}\partial_sF(f^\star(s),s)
\label{eq:obs-reduced-slow-flow}
\end{equation}
is the $\mathsf G_{\rm s}$ natural-gradient flow of $F_{\rm red}$, and along the
lifted curve $(f^\star(s(t)),s(t))$ the free energy decreases at the rate
$dF_{\rm red}\cdot\mathsf G_{\rm s}^{-1}dF_{\rm red}$. A $C^1$ branch exists near
$s_0\in\mathcal U_0$ whenever $F\in C^2$, $\partial^2_{ff}F(f^\star(s_0),s_0)$ is
positive definite, and $f^\star$ is continuous at $s_0$. If $\mathcal U_{\rm f}$ is
replaced by a compact set on which $F$ and $\partial_sF$ are continuous and the
conditional minimizer is unique, \eqref{eq:obs-envelope-identity} holds without
differentiability of $f^\star$. \status{ESTABLISHED}

\noindent\emph{Proof.}
At the minimizer $\partial_fF(f^\star(s),s)=0$, so the chain rule gives
$dF_{\rm red}=\partial_fF\cdot Df^\star+\partial_sF=\partial_sF$, and substitution in
\eqref{eq:obs-reduced-slow-flow} gives the flow and its rate. For the branch, the
implicit function theorem applied to $\partial_fF=0$ at $(f^\star(s_0),s_0)$ gives a
$C^1$ map whose graph is the whole critical set near that point, and continuity places
$f^\star(s)$ on it for $s$ near $s_0$. In the compact case, uniqueness and compactness
make $f^\star$ continuous, since a limit point $f'$ of $f^\star(s_n)$ with $s_n\to s$
satisfies $F(f',s)\leq F(f^\star(s),s)$. With $s'=s+\Delta s$, the difference
$F_{\rm red}(s')-F_{\rm red}(s)$ lies between
$F(f^\star(s'),s')-F(f^\star(s'),s)$ and $F(f^\star(s),s')-F(f^\star(s),s)$, and by the
mean value theorem and uniform continuity of $\partial_sF$ on a compact neighborhood
both equal $\partial_sF(f^\star(s),s)\Delta s+o(\lvert\Delta s\rvert)$; this is the
theorem of \citet[Ch.~III]{Danskin1967} in the unique-minimizer case. $\square$

\propositionheading{Rate-separated natural gradient and its reduced flow}{prop:obs-rate-separated-flow}
Let $F\in C^1(\mathcal U_{\rm f}\times\mathcal U_{\rm s})$, let
$\mathsf G_{\rm f}(f,s)$ and $\mathsf G_{\rm s}(s)$ be continuous symmetric
positive-definite block metrics with $\mathsf G_{\rm s}$ independent of $f$, and
for $\varepsilon>0$ set
$\mathsf G_\varepsilon=\varepsilon\mathsf G_{\rm f}\oplus\mathsf G_{\rm s}$. The
$\mathsf G_\varepsilon$ natural gradient of $F$ is the rate-separated flow
\begin{equation}
\varepsilon\dot f=-\mathsf G_{\rm f}^{-1}\partial_fF,
\qquad
\dot s=-\mathsf G_{\rm s}^{-1}\partial_sF,
\label{eq:obs-rate-separated-flow}
\end{equation}
which decreases $F$ at the rate
$\varepsilon^{-1}\partial_fF\cdot\mathsf G_{\rm f}^{-1}\partial_fF
+\partial_sF\cdot\mathsf G_{\rm s}^{-1}\partial_sF$ for every $\varepsilon>0$; its
orbit can depend on $\varepsilon$. Under the hypotheses of
\cref{lem:obs-envelope-slow-flow}, let
$\mathcal S^\star=\{(f^\star(s),s):s\in\mathcal U_0\}$. At every point of
$\mathcal S^\star$ and for every $\varepsilon>0$, the fast component of
\eqref{eq:obs-rate-separated-flow} vanishes and the slow component equals
$-\mathsf G_{\rm s}^{-1}dF_{\rm red}$. The reduced field
$\mathscr V^\star=(Df^\star v,v)$ with $v=-\mathsf G_{\rm s}^{-1}dF_{\rm red}$ is
tangent to $\mathcal S^\star$, and
\begin{equation}
d\operatorname{pr}_{\rm s}\mathscr V^\star
=-\operatorname{grad}^{\mathsf G_{\rm s}}F_{\rm red}\circ\operatorname{pr}_{\rm s}
\quad\text{on }\mathcal S^\star.
\label{eq:obs-reduced-semiconjugacy}
\end{equation}
Thus $(\operatorname{pr}_{\rm s}|_{\mathcal S^\star},1)$ is an oriented
semiconjugacy in the sense of \cref{def:hist-oriented-semiconjugacy}, with the
manifold $\mathcal S^\star$ in place of an open fine domain, and a conjugacy,
because $\operatorname{pr}_{\rm s}|_{\mathcal S^\star}$ is a diffeomorphism onto
$\mathcal U_0$. The speeds satisfy
$\lVert\mathscr V^\star\rVert^2_{\mathsf G_\varepsilon}
=\varepsilon\lVert Df^\star v\rVert^2_{\mathsf G_{\rm f}}+\lVert v\rVert^2_{\mathsf G_{\rm s}}$,
so the $\mathsf G_\varepsilon$-length of a reduced orbit is at least the
$\mathsf G_{\rm s}$-length of its image. These lengths decrease to the latter as
$\varepsilon\downarrow0$ on every compact time interval of the orbit, and on
any interval where the lifted length is finite at some positive
$\varepsilon_0$. The conjugacy and the length limit are properties of the lift
$\mathscr V^\star$ of the reduced field; they say nothing about orbits of
\eqref{eq:obs-rate-separated-flow} off $\mathcal S^\star$, whose relation to the
reduced orbit is \cref{open:obs-slow-manifold}. \status{ESTABLISHED}

\noindent\emph{Proof.}
Because $\mathsf G_\varepsilon^{-1}=\varepsilon^{-1}\mathsf G_{\rm f}^{-1}\oplus\mathsf G_{\rm s}^{-1}$,
the flow is \eqref{eq:obs-rate-separated-flow} with the stated dissipation, and the
orbit statement is \cref{prop:hist-scalar-mobility-orbit}, since
$\mathsf G_\varepsilon$ differs from $\mathsf G_1$ by an anisotropic rescaling. On
$\mathcal S^\star$, $\partial_fF=0$ and $\partial_sF=dF_{\rm red}$ by
\eqref{eq:obs-envelope-identity}, and $\mathsf G_{\rm s}$ does not depend on $f$. The
field $\mathscr V^\star$ is the differential of $s\mapsto(f^\star(s),s)$ applied to
$v$, hence tangent to $\mathcal S^\star$ with slow component $v$. The length
identity is the block form of $\mathsf G_\varepsilon$. For
$0<\varepsilon\le\varepsilon_0$ the lifted speed is bounded by its
$\varepsilon_0$ value, which is integrable under the stated hypothesis and
on compact time intervals by continuity. Dominated convergence gives the
length limit. $\square$

\remarkheading{The slow graph off the optimum and under correlated metrics}{rem:obs-rate-separation-witnesses}
Off $\mathcal S^\star$ the slow component $-\mathsf G_{\rm s}^{-1}\partial_sF(f,s)$
need not be a positive multiple of
$-\operatorname{grad}^{\mathsf G_{\rm s}}F_{\rm red}(s)$: for
$F=[(f-s_1)^2+s_1^2+s_2^2]/2$ with $\mathsf G_{\rm s}=I$,
$\partial_sF=(2s_1-f,s_2)$ and $\operatorname{grad}F_{\rm red}=(s_1,s_2)$ are not
parallel when $f\ne s_1$ and $s_2\ne0$, so the projection semiconjugates the
finite-$\varepsilon$ flow to the reduced flow only where that proportionality
holds. At points of $\mathcal S^\star$ the finite-$\varepsilon$ velocity $(0,v)$ is
tangent to $\mathcal S^\star$ only where $Df^\star v=0$, so $\mathcal S^\star$ is not
invariant in general. A correlated belief--model metric requires the rate
declaration to name its slow manifold. For $F(f,s)=[(f-s)^2+s^2]/2$,
whose conditional minimizer is $f=s$, and the metric $\mathsf G$ with unit diagonal
and off-diagonal entry $\mathrm c_2\in(0,1)$, the row-rated field $\varepsilon\dot f=-(\mathsf G^{-1}dF)_f$,
$\dot s=-(\mathsf G^{-1}dF)_s$ has, in the fast time $t/\varepsilon$, the critical set
\begin{equation}
\bigl\{(\mathsf G^{-1}dF)_f=0\bigr\}=\Bigl\{f=\frac{1+2\mathrm c_2}{1+\mathrm c_2}s\Bigr\},
\label{eq:obs-correlated-slow-manifold}
\end{equation}
which attracts at rate $1/(1-\mathrm c_2)$ and differs from the conditional
minimizer. This field need not descend $F$: at $\mathrm c_2=1/2$,
every $\varepsilon\in(0,7-4\sqrt3)$, and $(f,s)=(1+2\delta,1+\delta)$ with
$\delta=(1+\varepsilon)/4$, it gives
$dF/dt=(1+\varepsilon)^2/(12\varepsilon)-4/3>0$, so for those rates it
is the natural gradient of no metric. The congruence rescaling
$\mathsf G_\varepsilon=\Delta_\varepsilon\mathsf G\Delta_\varepsilon$ with
$\Delta_\varepsilon=\operatorname{diag}(\sqrt\varepsilon,1)$ gives a natural
gradient that descends $F$ for every $\varepsilon$, and the critical set of its
$\varepsilon=0$ fast subsystem is the conditional-minimizer set
$\{\partial_fF=0\}$. \status{ESTABLISHED}

\noindent\emph{Proof.}
The first two statements are direct differentiation. For the correlated metric,
$\mathsf G^{-1}$ has diagonal entries $1/(1-\mathrm c_2^2)$ and off-diagonal entries
$-\mathrm c_2/(1-\mathrm c_2^2)$, and $dF=(f-s,2s-f)$, so $(\mathsf G^{-1}dF)_f$ vanishes on the displayed line, and
its derivative in $f$ is $(1+\mathrm c_2)/(1-\mathrm c_2^2)=1/(1-\mathrm c_2)$.
Substituting the stated point into
$dF/dt=-\varepsilon^{-1}\partial_fF(\mathsf G^{-1}dF)_f-\partial_sF(\mathsf G^{-1}dF)_s$
at $\mathrm c_2=1/2$ gives the displayed value, which is positive exactly when
$\varepsilon^2-14\varepsilon+1>0$. Under congruence, $\mathsf G_\varepsilon$ is
positive definite for every $\varepsilon>0$, and in the fast time the fast
component of its natural gradient is
$-(1-\mathrm c_2^2)^{-1}(\partial_fF-\mathrm c_2\sqrt\varepsilon\partial_sF)$, whose
zero set at $\varepsilon=0$ is $\{\partial_fF=0\}$. $\square$

\propositionheading{Slow model means on the four-cycle and the model-link dichotomy}{prop:obs-slow-model-link-dichotomy}
In \cref{def:rg-formation-cycle-experiment} let the belief means $\bar k$ move at
rate one, the model means $\bar m$ at rate $\varepsilon_{\mathrm m}\in[0,1]$, and the
link translations $\check\omega$ at rate $\varepsilon_{\check\omega}\in[0,1]$: between clock times
\begin{equation}
d\mathbf x=-\bigl(\partial_{\bar k}U_S,\varepsilon_{\mathrm m}\partial_{\bar m}U_S,
\varepsilon_{\check\omega}\partial_{\check\omega} U_S\bigr)dt
+\sqrt{2\mathrm T}\bigl(dW^{\bar k},\varepsilon_{\mathrm m}^{1/2}dW^{\bar m},
\varepsilon_{\check\omega}^{1/2}dW^{\check\omega}\bigr)
\label{eq:obs-cycle-slow-model-kinetics}
\end{equation}
with one standard Brownian motion $W$ in $\R^{12}$, and the clock redraws the
assignment as before, leaving $\mathbf x$ unchanged. The rates are supplied, since
the Fisher metric of $\mathcal N(\bar m_i,\sigma^2_{\mathrm{ev}})$ in $\bar m_i$ is
$1/\sigma^2_{\mathrm{ev}}$ and a Fisher natural gradient would move the model means at
rate $1/2$. For $i\in\{0,2\}$ write $\bar k^\pm_i=(\bar k_{i+1}\pm\bar k_i)/\sqrt2$
and $\bar m^\pm_i=(\bar m_{i+1}\pm\bar m_i)/\sqrt2$, so that
$r^b_i=\sqrt2\bar k^-_i-\check\omega_i$ and $r^m_i=\sqrt2\bar m^-_i-\check\omega_i$.
\begin{enumerate}
\item[(i)] \emph{Held matching.} On the held assignment $S^{(0)}$ a Gaussian law
of $\mathbf x$ stays Gaussian and its mean follows
\eqref{eq:obs-cycle-slow-model-kinetics} without the noise; for
$\varepsilon_{\mathrm m},\varepsilon_{\check\omega}>0$, in the slow time
$\tau=\varepsilon_{\mathrm m}t$, this mean flow is \eqref{eq:obs-rate-separated-flow}
for $F=U_{S^{(0)}}$, which is the collective free energy of
\cref{prop:rg-formation-endogenous-cycle} at the Gaussian law up to terms
independent of the mean, with $\varepsilon=\varepsilon_{\mathrm m}$, $f=\bar k$,
$s=(\bar m,\check\omega)$, $\mathsf G_{\rm f}=I_4$ and
$\mathsf G_{\rm s}=I_4\oplus(\varepsilon_{\mathrm m}/\varepsilon_{\check\omega})I_4$. In the
orthonormal coordinates $(\bar k^-_i,\bar m^-_i,\check\omega_i)$ and
$(\bar k^+_i,\bar m^+_i)$ for $i\in\{0,2\}$, $\check\omega_1$, and $\check\omega_3$, the Hessian
of $U_{S^{(0)}}$ and the mobility are block diagonal, with the blocks
\begin{equation}
\begin{pmatrix}4+3\mathrm T&-2\mathrm T&-2\sqrt2\\-2\mathrm T&4+3\mathrm T&-2\sqrt2\\
-2\sqrt2&-2\sqrt2&4+\mathrm T\end{pmatrix},\qquad
\mathrm T\begin{pmatrix}3&-2\\-2&3\end{pmatrix},\qquad\mathrm T,
\label{eq:obs-cycle-held-blocks}
\end{equation}
so the mean flow is a direct sum of six flows and $r^b_i,r^m_i$ read only the block
of $e_i$. With $\varepsilon_{\check\omega}=\varepsilon_{\mathrm m}$ the split by mobility,
$f=\bar k$ and $s=(\bar m,\check\omega)$, is the only split for which the metric has the
form $\mathsf G_\varepsilon$ with blocks independent of $\varepsilon_{\mathrm m}$;
restricted to a selected-edge block it satisfies the three hypotheses of
\cref{open:obs-slow-manifold} with $\mathrm c_1=4$ for every $\mathrm T>0$, on every
closed ball whose interior contains the reduced orbit, while on the full sheet its
fast-block curvature along the pair means $\bar k^+_i$ is $3\mathrm T$.
\item[(ii)] \emph{Dichotomy at zero temperature.} Hold $S^{(0)}$, drop the
temperature terms, and start the mean flow
$\dot{\mathbf x}=-(\partial_{\bar k}V,\varepsilon_{\mathrm m}\partial_{\bar m}V,\varepsilon_{\check\omega}\partial_{\check\omega} V)$,
$V=V_{S^{(0)}}$, from $\mathbf x^\circ$, where every residual equals $-1$. Unless
$(\varepsilon_{\mathrm m},\varepsilon_{\check\omega})\in\{(0,0),(1,0)\}$, the flow matrix
of a selected-edge block has the simple eigenvalue $0$ on $(1,1,\sqrt2)$, which both
residual rows annihilate, and the two rates
\begin{equation}
\varsigma_\pm=2(1+\varepsilon_{\mathrm m}+\varepsilon_{\check\omega})
\pm2\sqrt{(1-\varepsilon_{\mathrm m})^2+\varepsilon_{\check\omega}^2},
\label{eq:obs-cycle-block-rates}
\end{equation}
and each internal residual is the combination of $e^{-\varsigma_-t}$ and
$e^{-\varsigma_+t}$ with value $-1$ and derivative
$4(\varepsilon_{\mathrm m}+\varepsilon_{\check\omega})$ (model) or $4(1+\varepsilon_{\check\omega})$
(belief) at $t=0$. With a fast link, $\varepsilon_{\check\omega}=1$, one has
$\varsigma_-\ge4-2\sqrt2$ for every $\varepsilon_{\mathrm m}\in[0,1]$ and
$|r^m_i(t)|\le e^{-(4-2\sqrt2)t}$: the link removes the model disagreement on the
belief clock even when the model means are frozen. With fast model means,
$\varepsilon_{\mathrm m}=1$ and $\varepsilon_{\check\omega}\in(0,1]$,
$r^m_i(t)=-e^{-(4+4\varepsilon_{\check\omega})t}$. With slow model means and links,
$\varepsilon_{\check\omega}=\varepsilon_{\mathrm m}\le1/2$, one has
$6\varepsilon_{\mathrm m}-2\varepsilon_{\mathrm m}^2\le\varsigma_-\le6\varepsilon_{\mathrm m}$,
$\varsigma_+\ge4$, and
\begin{equation}
|r^b_i(t)|\le(1+\varepsilon_{\mathrm m})e^{-4t}+\varepsilon_{\mathrm m},\qquad
\bigl(1-\tfrac32\varepsilon_{\mathrm m}\bigr)e^{-6\varepsilon_{\mathrm m}t}-\tfrac32\varepsilon_{\mathrm m}
\le|r^m_i(t)|\le e^{-(6\varepsilon_{\mathrm m}-2\varepsilon_{\mathrm m}^2)t}+\tfrac32\varepsilon_{\mathrm m},
\label{eq:obs-cycle-slow-link-bounds}
\end{equation}
so for $\varepsilon_{\mathrm m}\le1/25$ each matched pair has
$|r^b_i|\le2\varepsilon_{\mathrm m}$ and $|r^m_i|\ge1/2$ at every time from
$\tfrac14\log(4/\varepsilon_{\mathrm m})$ to $1/(12\varepsilon_{\mathrm m})$.
\item[(iii)] \emph{Persistence along every history.} Let $\mathrm T>0$, let
$\mathbf x_t$ be the sampled process from any initial configuration and
assignment, and let $\bar\varepsilon=\max(\varepsilon_{\mathrm m},\varepsilon_{\check\omega})>0$.
For an edge $e_i$ with $|r^m_i(\mathbf x_0)|\ge1$, numbers $\mathrm R\ge1$ and
$\mathrm R'>0$, and $0<t\le[4\bar\varepsilon(12\mathrm R+9\mathrm T\mathrm R')]^{-1}$,
the event that $|r^m_i(\mathbf x_u)|<1/2$ for some $u\le t$ while every residual
stays bounded by $\mathrm R$ and $|\mathbf x_u|\le\mathrm R'$ on $[0,t]$ has
probability at most $2\exp[-1/(192\mathrm T\bar\varepsilon t)]$, whatever the
assignment path, and at most $2e^{-1/(4\mathrm T)}$ at the largest such $t$.
\end{enumerate}
Among the three arms of (ii), a belief-agreeing, model-disagreeing interval of
order $1/\varepsilon_{\mathrm m}$ exists exactly in the arm
$\varepsilon_{\check\omega}=\varepsilon_{\mathrm m}\le1/2$: there it is the interval
of (ii), whose model disagreement (iii) protects along every assignment path up to a
time of order $1/\varepsilon_{\mathrm m}$ while residuals and configuration stay
bounded, and a fast link or fast model means remove the disagreement at a rate at
least $4-2\sqrt2$ whatever the other rate. Small rates alone do not suffice: for
$\varepsilon_{\mathrm m},\varepsilon_{\check\omega}\in[0,1/2]$,
\begin{equation}
4\varepsilon_{\mathrm m}+2\varepsilon_{\check\omega}-2\varepsilon_{\check\omega}^2
\le\varsigma_-\le4\varepsilon_{\mathrm m}+2\varepsilon_{\check\omega},
\qquad
|r^m_i(t)|\le\bigl(1+4(\varepsilon_{\mathrm m}+\varepsilon_{\check\omega})\bigr)
e^{-(4\varepsilon_{\mathrm m}+\varepsilon_{\check\omega})t},
\label{eq:obs-cycle-general-slow-rate}
\end{equation}
so $|r^m_i|\ge1/2$ at a time $\mathrm c_6/\varepsilon_{\mathrm m}$ with
$\mathrm c_6>0$ requires
$\varepsilon_{\check\omega}\le\varepsilon_{\mathrm m}\log(10)/\mathrm c_6$, and an
interval of order $1/\varepsilon_{\mathrm m}$ needs a link rate of order at most
$\varepsilon_{\mathrm m}$. \status{ESTABLISHED}

\noindent\emph{Proof.}
For (i), both endpoints of $e_0$ and of $e_2$ attend them under $S^{(0)}$ and no
receiver attends $e_1$ or $e_3$, so $V_{S^{(0)}}=2\mathcal D_{e_0}+2\mathcal D_{e_2}$
with $2\mathcal D_{e_i}=(\sqrt2\bar k^-_i-\check\omega_i)^2+(\sqrt2\bar m^-_i-\check\omega_i)^2$,
the reference term is invariant under the orthonormal change of coordinates, and
the evaluator term is
$\mathrm T\sum_{i\in\{0,2\}}[(\bar k^-_i-\bar m^-_i)^2+(\bar k^+_i-\bar m^+_i)^2]$;
collecting second derivatives gives \eqref{eq:obs-cycle-held-blocks}, and the
mobility stays diagonal because the new coordinates mix one channel at a time. On a
held assignment \eqref{eq:obs-cycle-slow-model-kinetics} is linear with constant
coefficients, so a Gaussian law stays Gaussian and its mean solves the noiseless
equation. Since $U_S(\mathbf x)=\frac12\mathbf x^\top\nabla^2U_S\mathbf x$, the
identity
$\Fenergy_{\mathrm T}(q)=\mathrm T\KL(q\Vert q^\star)+\mathrm B_{\mathrm T}-\mathrm T\log Z_{\mathrm{opt}}$
of \cref{prop:rg-formation-endogenous-cycle} at $q=\mathcal N(\mu,\Sigma)\otimes\delta_S$
gives $\Fenergy_{\mathrm T}=U_S(\mu)+\frac12\Tr(\nabla^2U_S\Sigma)-\frac{\mathrm T}2\log\det\Sigma$
up to a constant. Dividing the mean equation by $\varepsilon_{\mathrm m}$ gives the
slow-time form with the metric
$\varepsilon_{\mathrm m}I_4\oplus I_4\oplus(\varepsilon_{\mathrm m}/\varepsilon_{\check\omega})I_4$,
whose blocks are independent of $\varepsilon_{\mathrm m}$ exactly when $f$ consists
of the unit-rate coordinates. On a selected-edge block with
$\varepsilon_{\check\omega}=\varepsilon_{\mathrm m}$, $F$ is the quadratic form of the first
block of \eqref{eq:obs-cycle-held-blocks}, positive definite for $\mathrm T>0$ since
its leading minors are $4+3\mathrm T$, $(4+\mathrm T)(4+5\mathrm T)$ and
$\mathrm T(8+\mathrm T)(4+5\mathrm T)$; $\partial^2_{ff}F=4+3\mathrm T\ge4\mathsf G_{\rm f}$;
the reduced matrix is a Schur complement, hence positive definite, so $|s|$ is
nonincreasing along the reduced flow and every closed ball whose interior contains
the initial slow state serves as $\mathcal K$; and the frozen fast flow is linear
with coefficient $4+3\mathrm T$, so every fast state lies in its basin. The
pair-mean block gives the curvature $3\mathrm T$ along $\bar k^+_i$.

For (ii), the block flow matrix is $\operatorname{diag}(1,\varepsilon_{\mathrm m},\varepsilon_{\check\omega})$
times the first block of \eqref{eq:obs-cycle-held-blocks} at $\mathrm T=0$, whose
kernel is spanned by $(1,1,\sqrt2)$, on which $r^b_i=r^m_i=0$. Expanding the
determinant, its characteristic polynomial is
$\varsigma[\varsigma^2-4(1+\varepsilon_{\mathrm m}+\varepsilon_{\check\omega})\varsigma+8(2\varepsilon_{\mathrm m}+\varepsilon_{\check\omega}+\varepsilon_{\mathrm m}\varepsilon_{\check\omega})]$,
whose quadratic factor has discriminant
$16[(1-\varepsilon_{\mathrm m})^2+\varepsilon_{\check\omega}^2]$ and constant term vanishing
only at $(0,0)$, so its roots are \eqref{eq:obs-cycle-block-rates} and the
eigenvalues are distinct outside the two excluded pairs. The matrix is then
diagonalizable and the residual rows annihilate the kernel, so each internal residual
is $\mathrm c_2e^{-\varsigma_-t}+\mathrm c_3e^{-\varsigma_+t}$ with
$\mathrm c_2+\mathrm c_3=-1$ and $\varsigma_-\mathrm c_2+\varsigma_+\mathrm c_3$ equal
to minus the derivative at $t=0$, which is
$a\cdot\operatorname{diag}(1,\varepsilon_{\mathrm m},\varepsilon_{\check\omega})(-2\sqrt2,-2\sqrt2,4)^\top$
for the residual row $a$; hence $\mathrm c_3=(\varsigma_--\mathrm c_5)/(\varsigma_+-\varsigma_-)$
with $\mathrm c_5=4(\varepsilon_{\mathrm m}+\varepsilon_{\check\omega})$ for $r^m_i$ and
$\mathrm c_5=4(1+\varepsilon_{\check\omega})$ for $r^b_i$. For $\varepsilon_{\check\omega}=1$ put
$\mathrm c_4=\sqrt{(1-\varepsilon_{\mathrm m})^2+1}$, so
$\varsigma_\pm=4+2\varepsilon_{\mathrm m}\pm2\mathrm c_4$; $\varsigma_-$ increases in
$\varepsilon_{\mathrm m}$ from $4-2\sqrt2$, and
$\varsigma_-\le4+4\varepsilon_{\mathrm m}\le\varsigma_+$ because
$\mathrm c_4\ge\varepsilon_{\mathrm m}$, so both coefficients of $r^m_i$ are
nonpositive with sum $-1$. For $\varepsilon_{\mathrm m}=1$,
$\varsigma_-=4=\mathrm c_5-4\varepsilon_{\check\omega}$ gives $\mathrm c_3=-1$. For
$\varepsilon_{\check\omega}=\varepsilon_{\mathrm m}\le1/2$ put
$\mathrm c_4=\sqrt{(1-\varepsilon_{\mathrm m})^2+\varepsilon_{\mathrm m}^2}$, so that
$1-\varepsilon_{\mathrm m}\le\mathrm c_4\le1-\varepsilon_{\mathrm m}+\varepsilon_{\mathrm m}^2$
by squaring, $\varsigma_\pm=2+4\varepsilon_{\mathrm m}\pm2\mathrm c_4$ and
$\varsigma_+-\varsigma_-=4\mathrm c_4\ge4(1-\varepsilon_{\mathrm m})$, which gives the
stated brackets; the coefficient of $e^{-\varsigma_+t}$ in $r^m_i$ is
$(\varsigma_--8\varepsilon_{\mathrm m})/(4\mathrm c_4)$, of absolute value at most
$\varepsilon_{\mathrm m}(1+\varepsilon_{\mathrm m})/(2(1-\varepsilon_{\mathrm m}))\le\frac32\varepsilon_{\mathrm m}$,
and the coefficient of $e^{-\varsigma_-t}$ in $r^b_i$ is
$(\varsigma_--4\varepsilon_{\mathrm m})/(4\mathrm c_4)\in[0,\varepsilon_{\mathrm m}]$.
Since each pair of coefficients sums to $-1$, \eqref{eq:obs-cycle-slow-link-bounds}
follows, and the last claim is \eqref{eq:obs-cycle-slow-link-bounds} at the two
times, using $(1+\varepsilon_{\mathrm m})/4\le1$ and
$e^{-1/2}(1-\frac3{50})-\frac3{50}>\frac12$.

For (iii), write $a$ for the row of $r^m_i$, with entries $-1$ and $1$ at
$\bar m_i$ and $\bar m_{i+1}$ and $-1$ at $\check\omega_i$. Since the clock leaves
$\mathbf x$ unchanged, \eqref{eq:obs-cycle-slow-model-kinetics} gives, along every
assignment path,
$r^m_i(\mathbf x_t)-r^m_i(\mathbf x_0)=-\int_0^ta\cdot(\partial_{\bar k}U_{S_u},\varepsilon_{\mathrm m}\partial_{\bar m}U_{S_u},\varepsilon_{\check\omega}\partial_{\check\omega} U_{S_u})(\mathbf x_u)du+\sqrt{2\mathrm T(2\varepsilon_{\mathrm m}+\varepsilon_{\check\omega})}B_t$
with $B$ a standard Brownian motion. With $\mathrm n_j(S)\in\{0,1,2\}$ the number of
receivers attending $e_j$ under $S$,
$\partial_{\bar m_j}U_S=\mathrm n_{j-1}(S)r^m_{j-1}-\mathrm n_j(S)r^m_j+3\mathrm T\bar m_j-2\mathrm T\bar k_j$
and $\partial_{\check\omega_j}U_S=-\mathrm n_j(S)(r^b_j+r^m_j)+\mathrm T\check\omega_j$, so on the
stated event the integrand is bounded by
$\varepsilon_{\mathrm m}(8\mathrm R+8\mathrm T\mathrm R')+\varepsilon_{\check\omega}(4\mathrm R+\mathrm T\mathrm R')\le\bar\varepsilon(12\mathrm R+9\mathrm T\mathrm R')$,
using $3|\bar m_j|+2|\bar k_j|\le\sqrt{13}|\mathbf x|$. For $t$ as stated the drift
moves the residual by at most $1/4$, so a crossing below $1/2$ requires
$\sup_{u\le t}|B_u|\ge[4\sqrt{2\mathrm T(2\varepsilon_{\mathrm m}+\varepsilon_{\check\omega})}]^{-1}\ge[4\sqrt{6\mathrm T\bar\varepsilon}]^{-1}$.
Doob's maximal inequality applied to the martingale $\exp(\theta B_u-\theta^2u/2)$
with $\theta=\mathrm y/t$ gives $\Pr(\sup_{u\le t}B_u\ge\mathrm y)\le e^{-\mathrm y^2/(2t)}$,
and with $\mathrm y^2=1/(96\mathrm T\bar\varepsilon)$ the two-sided bound is
$2e^{-1/(192\mathrm T\bar\varepsilon t)}$; at the largest $t$ the exponent is
$-(12\mathrm R+9\mathrm T\mathrm R')/(48\mathrm T)\le-1/(4\mathrm T)$. The first
closing sentence collects (ii) and (iii), the fast-means case having rate
$4+4\varepsilon_{\check\omega}\ge4$. For \eqref{eq:obs-cycle-general-slow-rate} put
$\mathrm c_4=\sqrt{(1-\varepsilon_{\mathrm m})^2+\varepsilon_{\check\omega}^2}$; since
$2(1-\varepsilon_{\mathrm m})\ge1$, squaring gives
$1-\varepsilon_{\mathrm m}\le\mathrm c_4\le1-\varepsilon_{\mathrm m}+\varepsilon_{\check\omega}^2$,
and $\varsigma_-=2(1+\varepsilon_{\mathrm m}+\varepsilon_{\check\omega})-2\mathrm c_4$
gives the bracket, whose lower end is at least
$4\varepsilon_{\mathrm m}+\varepsilon_{\check\omega}$ because
$\varepsilon_{\check\omega}\le1/2$. At $(0,0)$ the residual is constant. Otherwise
the coefficient of $e^{-\varsigma_+t}$ in $r^m_i$ is
$\mathrm c_3=(\varsigma_--4\varepsilon_{\mathrm m}-4\varepsilon_{\check\omega})/(4\mathrm c_4)$,
and $0\le\varsigma_-\le4\varepsilon_{\mathrm m}+2\varepsilon_{\check\omega}$ with
$\mathrm c_4\ge1/2$ give
$|\mathrm c_3|\le2(\varepsilon_{\mathrm m}+\varepsilon_{\check\omega})$; the other
coefficient is $-1-\mathrm c_3$, so
$|r^m_i(t)|\le(1+2|\mathrm c_3|)e^{-\varsigma_-t}$, which is the second bound. At
$t=\mathrm c_6/\varepsilon_{\mathrm m}$ that bound is at most
$5e^{-\mathrm c_6\varepsilon_{\check\omega}/\varepsilon_{\mathrm m}}$, which is below
$1/2$ when $\varepsilon_{\check\omega}>\varepsilon_{\mathrm m}\log(10)/\mathrm c_6$.
$\square$

In the latent reading of \cref{ch:links}, the link translations are generators of the translation
instance with anchor zero and no loop term, and the supplied rate $\varepsilon_{\check\omega}$ places
their recognition means in the slow block on the held matching sheet; the natural-gradient recognition
flow of the generators assigns belief means, model means, and link means one common rate on an invariant
slice, so the separation used here comes from the supplied mobilities (\cref{rem:link-rates}).

\openproblemheading{Slow-manifold reduction of the collective free energy}{open:obs-slow-manifold}
Consider the rate-separated flow \eqref{eq:obs-rate-separated-flow} under three
hypotheses. First, $F\in C^2$ and the block metric
$\mathsf G_{\rm f}(f,s)\oplus\mathsf G_{\rm s}(s)$ has $C^1$ symmetric
positive-definite blocks with $\mathsf G_{\rm s}$ independent of $f$. Second, there
are a compact set $\mathcal K\subset\mathcal U_0$, a $C^1$ branch $f^\star$ of
conditional minimizers, and $\mathrm c_1>0$ with
$\partial^2_{ff}F(f^\star(s),s)\succeq\mathrm c_1\mathsf G_{\rm f}(f^\star(s),s)$ for
all $s\in\mathcal K$. Third, the reduced flow \eqref{eq:obs-reduced-slow-flow} from
$s_0$ stays in the interior of $\mathcal K$ on $[0,t_1]$, and the initial fast state
lies in the basin of $f^\star(s_0)$ for the frozen fast flow
$\dot f=-\mathsf G_{\rm f}^{-1}\partial_fF(f,s_0)$. Under the second hypothesis the
linearized frozen fast flow has spectrum in $(-\infty,-\mathrm c_1]$, so $f^\star(s)$
is isolated and uniformly asymptotically stable on $\mathcal K$, and Tikhonov's
theorem gives $s_\varepsilon(t)\to s(t)$ for $0\leq t\leq t_1$ and
$f_\varepsilon(t)\to f^\star(s(t))$ for $0<t\leq t_1$ as $\varepsilon\to0$
\citep{Tikhonov1952}; for smooth data Fenichel's theory gives a locally invariant
slow manifold $O(\varepsilon)$-close to a compact normally hyperbolic piece of
$\mathcal S^\star$ \citep{fenichel1979geometric}. Their application at $c_*$ is open
beyond the held matching sheets of \cref{prop:obs-slow-model-link-dichotomy}(i),
whose selected-edge blocks meet the three hypotheses with a curvature constant
uniform in the temperature, so that Tikhonov's theorem applies to them, while the
split by mobility on the full sheet has a constant of order $\mathrm T$:
\cref{prop:obs-attention-reduced-hessian} shows that the second hypothesis can fail
after the source rows are optimized even when every source energy and every
frozen-weight curvature is positive, the sampled dynamics with source jumps of
\cref{ch:finite-depth-formation} lie outside this deterministic statement, and no
explicit bound on the slow orbit or its Fisher duration in terms of $\varepsilon$ is
available. It is settled in the class of the sampled dynamics of
\cref{def:rg-formation-cycle-experiment} under the mobility of
\cref{prop:obs-slow-model-link-dichotomy} with $\varepsilon_{\check\omega}=\varepsilon_{\mathrm m}$
and no parameter learning, where the split by mobility, $f=\bar k$ and
$s=(\bar m,\check\omega)$, and \eqref{eq:obs-reduced-slow-flow} for $F=U_S$ under the
current assignment $S$ define the reduced slow flow of an assignment path: positively,
by a bound on the distance between the slow coordinates of the sampled process and
that reduced flow over a bounded slow-time interval, explicit in
$\varepsilon_{\mathrm m}$, $\mathrm T$, and the number of assignment changes, tending
to zero with $\varepsilon_{\mathrm m}$ and $\mathrm T$, and holding, for every
positive integer $\mathrm n$, with probability tending to one on the event that the
assignment path has at most $\mathrm n$ changes per unit slow time; negatively, by
a positive integer $\mathrm n$, a sequence of mobility ratios and temperatures
tending to zero, and initial states from which, with probability bounded below, the
assignment path has at most $\mathrm n$ changes per unit slow time and that
distance stays bounded below over a bounded slow-time interval. The
learning case, in which the generative parameter is the slow coordinate and the
reduced object is the recognition law at the frozen parameter, lies outside this
class.

\section{Parent-sensitive inference and the prepared parent}
\label{sec:rg-memory-parent-mechanism}

\definitionheading{Parent-sensitive sweep}{def:rg-memory-parent-sweep}
Let $\Xi$ be a nonempty finite set of fine states with a parent map
$\operatorname{par}:\Xi\to\mathsf Z$ that retains the complete decorated parent,
keeping only parent values with nonempty fibers. Take a prior $p^0$ on $\Xi$
with $p^0(\xi)>0$, normalized conditional laws $\psi_\xi$ of an unobserved record
$r$ in a finite set $\mathsf O'$, and record likelihoods $L_o(\xi)>0$ at the
observation, none of which reads recognition. The generative law of $(\xi,r,o)$
has density $p^0(\xi)\psi_\xi(r)L_o(\xi)$, and its evidence and posterior are
\begin{equation}
\mathcal Z_o=\sum_\xi p^0(\xi)L_o(\xi),
\qquad
p(\xi)=\frac{p^0(\xi)L_o(\xi)}{\mathcal Z_o}.
\label{eq:rg-memory-target}
\end{equation}
On an undirected graph of allowed moves take symmetric proposal weights
$\mathsf a_{\xi\xi'}=\mathsf a_{\xi'\xi}\geq0$ with off-diagonal row sums at most
$1/2$. The Metropolis kernel, the parent projection, the posterior lift, the
refresh, and the sweep are
\begin{equation}
\begin{aligned}
\mathsf M_{\xi\xi'}&=\mathsf a_{\xi\xi'}\min\bigl(1,p(\xi')/p(\xi)\bigr)\quad(\xi\ne\xi'),
&\mathsf M_{\xi\xi}&=1-\sum_{\xi'\ne\xi}\mathsf M_{\xi\xi'},\\
\mathsf C_{\xi z}&=\mathbf 1_{\operatorname{par}(\xi)=z},
&\mathsf J_{z\xi}&=\frac{p(\xi)\mathbf 1_{\operatorname{par}(\xi)=z}}{p_Z(z)},
\qquad p_Z=p\mathsf C,\\
\mathsf D&=\mathsf C\mathsf J,
&\mathsf T&=\mathsf M\mathsf D.
\end{aligned}
\label{eq:rg-memory-sweep}
\end{equation}
The sweep depends on the parent through the conditional laws $p(\cdot\given z)$. The
marked version of a kernel $\mathsf K$ on
$\Xi$ is $\widetilde{\mathsf K}((\xi,r),(\xi',r'))=\mathsf K_{\xi\xi'}\psi_{\xi'}(r')$.
The proposal and the update order are declared choices.

\hypothesisheading{Instantaneous within-parent refresh}{hyp:rg-memory-instantaneous-refresh}
In this section the refresh is the exact kernel $\mathsf D=\mathsf C\mathsf J$:
each application resamples the fine state from its exact conditional law given
the current complete parent. This is the adiabatic elimination of the
within-parent fine state, whose timescale is infinitely shorter than the parent's.
It excludes finite
within-fiber mixing, for example the partial refresh
$\mathsf D_\omega=(1-\omega)I+\omega\mathsf C\mathsf J$ with $0\leq\omega<1$, which
is equally normalized, $p$-reversible, and parent preserving.

\theoremheading{Stationarity, exact descent, and finite-component convergence}{thm:rg-memory-inference-descent}
\begin{enumerate}
\item[(i)] For a recognition law $q(\xi)\psi'_\xi(r)$ the free energy is
\begin{equation}
\Fenergy(q\otimes\psi';o)=\KL(q\Vert p)+\E_q\KL(\psi'_\xi\Vert\psi_\xi)
-\log\mathcal Z_o,
\label{eq:rg-memory-fixed-vfe}
\end{equation}
so matching the record conditionals, $\psi'=\psi$, is a specialization of the
joint objective and not the removal of its dependence terms.
\item[(ii)] The kernels $\mathsf M$, $\mathsf D$, $\mathsf T$ are normalized and
preserve $p$, and $\mathsf M$ and $\mathsf D$ are $p$-reversible; their marked
versions are normalized and preserve $p(\xi)\psi_\xi(r)$.
\item[(iii)] $q\mathsf D$ minimizes $\KL(\cdot\Vert p)$ among laws with parent
marginal $q\mathsf C$, with $\KL(q\mathsf D\Vert p)=\KL(q\mathsf C\Vert p_Z)$.
\item[(iv)] Call a component a connected component of the union of the positive
off-diagonal edges of $\mathsf M$ and $\mathsf D$. For every initial $q$, $\mathsf T$
preserves component masses, and for every $p$-stationary kernel $\mathsf K$ on
$\Xi$, with $q'=q\mathsf K$ and reversal
$\mathsf K^*_{\xi'\xi}=p(\xi)\mathsf K_{\xi\xi'}/p(\xi')$,
\begin{equation}
\KL(q\Vert p)-\KL(q\mathsf K\Vert p)
=\sum_{\xi':q'(\xi')>0}q'(\xi')
\KL\bigl(q(\cdot\given\xi')\Vert\mathsf K^*_{\xi'\cdot}\bigr),
\label{eq:rg-memory-reverse-conditional-loss}
\end{equation}
where $q(\xi\given\xi')=q(\xi)\mathsf K_{\xi\xi'}/q'(\xi')$. Hence the marked
sweep maps $q\otimes\psi'$ to $(q\mathsf T)\otimes\psi$ and lowers
\eqref{eq:rg-memory-fixed-vfe} by
$\KL(q\Vert p)-\KL(q\mathsf T\Vert p)+\E_q\KL(\psi'_\xi\Vert\psi_\xi)\geq0$,
which vanishes exactly when $q(\xi)=q(\Xi')p(\xi\given\Xi')$ on each component
$\Xi'$ and $\psi'_\xi=\psi_\xi$ for every $\xi$ with $q(\xi)>0$. For matched
record conditionals, $\psi'=\psi$, this is the statement that $q\mapsto q\mathsf T$
decreases $\KL(q\Vert p)$, strictly unless $q$ has that form.
\item[(v)] On each component, $q\mathsf T^n$ converges geometrically to that
conditional target.
\end{enumerate}
\status{ESTABLISHED}

\noindent\emph{Proof.}
For (i), the posterior of $(\xi,r)$ given $o$ is $p(\xi)\psi_\xi(r)$ because $r$
is unobserved and conditionally independent of $o$, and the evidence is
$\mathcal Z_o$; \cref{def:elbo-extended} and the chain rule
\cref{lem:prob-kl-chain-rule} give the display. For (ii),
$p(\xi)\mathsf M_{\xi\xi'}=\mathsf a_{\xi\xi'}\min(p(\xi),p(\xi'))=p(\xi')\mathsf M_{\xi'\xi}$;
$\mathsf J\mathsf C=I$ gives $p\mathsf D=p$, and
$p(\xi)\mathsf D_{\xi\xi'}=p(\xi)p(\xi')/p_Z(z)$ on a common parent fiber $z$.
Stationarity survives composition; summing the terminal record of a marked kernel
proves normalization and summing the initial record proves invariance. Part (iii) is
\cref{prop:elbo-exact-e-coordinate} for the block of fine coordinates given the
parent: $q\mathsf D$ has parent marginal $q\mathsf C\mathsf J\mathsf C=q\mathsf C$ and
conditionals $p(\cdot\given z)$, and the chain rule over the parent gives the value.
For (iv), the input--output laws $q(\xi)\mathsf K_{\xi\xi'}$ and
$p(\xi)\mathsf K_{\xi\xi'}$ have divergence $\KL(q\Vert p)$, and disintegrating them
over the output with \cref{lem:prob-kl-chain-rule} gives
\eqref{eq:rg-memory-reverse-conditional-loss}. With positive holding, equality forces
$q/p$ to be constant across every positive edge; equality for $\mathsf M$ makes $q$
stationary under $\mathsf M$, and equality for the subsequent $\mathsf D$ makes the
ratio constant on parent fibers, which is constancy on components. The marked sweep
sends $q(\xi)\psi'_\xi(r)$ to
$\sum_{\xi,r}q(\xi)\psi'_\xi(r)\mathsf T_{\xi\xi'}\psi_{\xi'}(r')=(q\mathsf T)(\xi')\psi_{\xi'}(r')$,
so by (i) the free energy before minus after is the displayed sum, whose two
terms are nonnegative, the first by \eqref{eq:rg-memory-reverse-conditional-loss}
and the second as an average of relative entropies, and which vanishes exactly
when both terms do. For (v), positive
holding makes every $\mathsf M$ or $\mathsf D$ edge a positive $\mathsf T$ edge, so on
a component $\Xi'$ of $n$ states $\mathsf T^k>0$ for $k=\max(1,n-1)$. With target
$p_{\Xi'}=p(\cdot\given\Xi')$ and
$\mathrm c_1=\min_{\xi\xi'}(\mathsf T^k)_{\xi\xi'}/p_{\Xi'}(\xi')\in(0,1]$, one has
$\mathsf T^k=\mathrm c_1\mathbf 1p_{\Xi'}+(1-\mathrm c_1)\mathsf K'$ on $\Xi'$ with a
$p_{\Xi'}$-stationary stochastic $\mathsf K'$ (or $\mathsf T^k=\mathbf 1p_{\Xi'}$ when
$\mathrm c_1=1$), and convexity with data processing give
$\KL(q\mathsf T^k\Vert p_{\Xi'})\leq(1-\mathrm c_1)\KL(q\Vert p_{\Xi'})$ for laws $q$ on
$\Xi'$. $\square$

\theoremheading{The prepared parent is Markov}{thm:rg-memory-prepared-parent}
Under \cref{hyp:rg-memory-instantaneous-refresh}, the reduced sweep
$\overline{\mathsf T}:=\mathsf J\mathsf M\mathsf C$ is reversible under $p_Z$ and
\begin{equation}
\mathsf J\mathsf T=\overline{\mathsf T}\mathsf J,
\qquad
(\nu\mathsf J)\mathsf T^n=\nu\overline{\mathsf T}^n\mathsf J.
\label{eq:rg-memory-parent-semiconjugacy}
\end{equation}
For a start $q_0=\nu\mathsf J$ the parent history $Z_t=\operatorname{par}(\xi_t)$
of the chain with kernel $\mathsf T$ is Markov with transition $\overline{\mathsf T}$,
and for an arbitrary start it is so after the first completed refresh.
\status{ESTABLISHED}

\noindent\emph{Proof.}
With $\mathsf T=\mathsf M\mathsf C\mathsf J$ and $\mathsf J\mathsf C=I$,
$\mathsf J\mathsf T=\mathsf J\mathsf M\mathsf C\mathsf J=\overline{\mathsf T}\mathsf J$,
and induction gives the second identity. Moreover $p_Z(z)\overline{\mathsf T}_{zz'}$
is the sum of $p(\xi)\mathsf M_{\xi\xi'}$ over $\operatorname{par}(\xi)=z$ and
$\operatorname{par}(\xi')=z'$, which is symmetric by reversibility of $\mathsf M$. The
final refresh makes the fine law given the whole parent history $\mathsf J_{z\cdot}$,
with $z$ the last parent, so the next parent row is
$\mathsf J_{z\cdot}\mathsf M\mathsf C=\overline{\mathsf T}_{z\cdot}$, and induction over
histories gives the Markov property. $\square$

The refresh holds the fast fine state at its conditional optimum given the slow parent
(\cref{thm:rg-memory-inference-descent}(iii)), so $\mathsf J\mathsf T=\overline{\mathsf T}\mathsf J$
is the discrete instantaneous-equilibration case of \eqref{eq:obs-reduced-semiconjugacy}.

\remarkheading{Finite refresh leaves parent memory}{rem:rg-memory-partial-refresh}
\Cref{hyp:rg-memory-instantaneous-refresh} cannot be dropped. Take
$\Xi=\{1,2,3\}$ with $\operatorname{par}(1)=\operatorname{par}(2)=z_1$ and
$\operatorname{par}(3)=z_2$, uniform $p$, the single proposal
$\mathsf a_{23}=1/2$, the partial sweep
$\mathsf T_\omega=\mathsf M\mathsf D_\omega$, and a start $\nu\mathsf J$ with
$\nu(z_1),\nu(z_2)>0$. Then
\begin{equation}
\mathbb P(Z_2=z_2\given Z_1=z_1,Z_0=z_1)=1/6+\omega/12,
\qquad
\mathbb P(Z_2=z_2\given Z_1=z_1,Z_0=z_2)=1/2-\omega/4.
\label{eq:rg-memory-partial-refresh-witness}
\end{equation}
The dependence on $Z_0$ is $(1-\omega)/3$, equal to $7/30$ at $\omega=3/10$, and it
vanishes only at $\omega=1$, where both values equal
$\overline{\mathsf T}_{z_1z_2}=1/4$. For $\omega<1$ the parent history is not
Markov. It remains a deterministic function of the fine Markov chain with kernel
$\mathsf T_\omega$, so the hidden-history law of \cref{sec:rg-memory-hidden-history}
retains it exactly, and the memory cost of a first-order parent closure is the
defect $D_N^{(1)}$ of \cref{thm:rg-memory-projection}, whose value
$I(Z_2;Z_0\given Z_1)$ at $N=3$ is positive. \status{ESTABLISHED}

\noindent\emph{Proof.}
Only state $2$ moves under $\mathsf M$, to $3$ with probability $1/2$, and
$\mathsf D_\omega$ keeps a state of $\{1,2\}$ with probability $1-\omega/2$. Given
$Z_1=z_1$, the state before the second sweep is $2$ with probability
$1-\omega/2$ if $Z_0=z_2$, since it moved from $3$, and with probability
$1/3+\omega/6$ if $Z_0=z_1$, since after $\mathsf M$ it is $2$ with probability
$1/3$ and the refresh redraws it uniformly with probability $\omega$; multiplying
by $1/2$ gives
\eqref{eq:rg-memory-partial-refresh-witness}. Both conditioning events have positive
probability, so $I(Z_2;Z_0\given Z_1)>0$. $\square$

\section*{Physicist's summary}

Fast and slow are declared, not derived. The collective free energy is then a
potential for every positive-definite mobility. The reduced gradient flow is
conjugate to its tangent lift on the graph of the conditional optimum; the
finite-rate flow shares its slow component there but can leave that graph.
Instantaneous within-parent refresh gives the Markov parent proved here.
On the four-cycle, slow model means lag the beliefs only if the links
are slow too; a fast link carries the model agreement on the belief clock.
